% Part: sets-functions-relations
% Chapter: relations
% Section: special-properties

\documentclass[../../../include/open-logic-section]{subfiles}

\begin{document}

\olfileid[tr]{sfr}{rel}{prp}
\olsection{Bağıntıların Özel Özellikleri}

\begin{intro}
Bazı bağıntı türlerine öylesine sık rastlanır ki bunlara özel adlar verilmiştir.
Örneğin $\le$ ve~$\subseteq$, kendi tanım kümelerindeki elemanları (söz gelimi
$\le$ için $\Nat$'ın, $\subseteq$ için $\Pow{A}$'nın elemanlarını) benzer
biçimlerde birbirine bağlar. Bu bağıntıların tam olarak hangi bakımlardan
benzeştiğini ve hangi bakımlardan ayrıldığını görmek için onları, bağıntıların
taşıyabileceği bazı özel özelliklere göre sınıflandırırız. Bu özel özelliklerin
bazılarının ve bunların birleşimlerinin özellikle önemli olduğu ortaya çıkar:
sıralamalar ve denklik bağıntıları.
\end{intro}

\begin{defn}[Yansımalılık]
Bir $R \subseteq A^2$ bağıntısı, ancak ve ancak her $x \in A$ için $Rxx$
geçerliyse \emph{yansımalıdır}.
\end{defn}

\begin{defn}[Geçişlilik]
Bir $R \subseteq A^2$ bağıntısı, ancak ve ancak $Rxy$ ile $Ryz$ birlikte geçerli
olduğunda $Rxz$ de geçerliyse \emph{geçişlidir}.
\end{defn}

\begin{defn}[Simetriklik]
Bir $R \subseteq A^2$ bağıntısı, ancak ve ancak $Rxy$ geçerli olduğunda
$Ryx$ de geçerliyse \emph{simetriktir}.
\end{defn}

\begin{defn}[Ters simetriklik]
Bir $R \subseteq A^2$ bağıntısı, ancak ve ancak hem $Rxy$ hem de $Ryx$ geçerli
olduğunda $x=y$ ise \emph{ters simetriktir} (başka bir deyişle,
$x\neq y$ ise ya $\lnot Rxy$ ya da $\lnot Ryx$ olur).
\end{defn}

\begin{explain}
Simetrik bir bağıntıda $Rxy$ ile $Ryx$ ya her zaman birlikte geçerlidir ya da
ikisi de geçerli değildir. Ters simetrik bir bağıntıda $Rxy$ ile $Ryx$'in
birlikte geçerli olmasının tek yolu $x=y$ olmasıdır. Bunun, $x=y$ olduğunda
$Rxy$ ile $Ryx$'in geçerli olmasını \emph{gerektirmediğine}; yalnızca bunu
dışlamadığına dikkat edin. Dolayısıyla ters simetrik bir bağıntı yansımalı
olabilir, ancak her ters simetrik bağıntı yansımalı değildir. Ters simetrik
olmakla yalnızca simetrik olmamak da farklı koşullardır. Nitekim bir bağıntı
aynı anda hem simetrik hem de ters simetrik olabilir (örneğin özdeşlik bağıntısı
böyledir).
\end{explain}

\begin{defn}[Bağlantılılık]
Bir $R \subseteq A^2$ bağıntısı, bütün $x,y\in A$ için $x\neq y$ olduğunda
$Rxy$ ya da~$Ryx$ ifadelerinden en az biri geçerliyse \emph{bağlantılıdır}.
\end{defn}

\begin{prob}
Şu özellikleri taşıyan bağıntılara örnekler verin: (a) yansımalı ve simetrik
olup geçişli olmayan, (b) yansımalı ve ters simetrik olan, (c) ters simetrik ve
geçişli olup yansımalı olmayan ve (d) yansımalı, simetrik ve geçişli olan.
Sayılar veya kümeler üzerindeki bağıntıları kullanmayın.
\end{prob}

\begin{defn}[Yansımasızlık]
Bir $R \subseteq A^2$ bağıntısına, bütün $x \in A$ için $Rxx$ geçerli değilse
\emph{yansımasız} denir.
\end{defn}

\begin{defn}[Asimetriklik]
Bir $R \subseteq A^2$ bağıntısına, hiçbir $x,y\in A$ ikilisi için $Rxy$ ile
$Ryx$ birlikte geçerli olmuyorsa \emph{asimetrik} denir.
\end{defn}

$A \neq \emptyset$ ise $A$ üzerindeki hiçbir yansımasız bağıntının yansımalı
olmadığına ve $A$ üzerindeki her asimetrik bağıntının aynı zamanda ters
simetrik olduğuna dikkat edin. Bununla birlikte yansımalı da yansımasız da
olmayan $R \subseteq A^2$ bağıntıları ve asimetrik olmayan ters
simetrik bağıntılar vardır.

\end{document}
